% Sección 3 (Rol 2: Diego, pasos 1-4; Amelie, pasos 5-6). Reemplaza el cuerpo de
% \section{Análisis exploratorio} en main.tex (la figura y el marcador del párrafo).
% Párrafo condensado de los textos ya aprobados en informe/reportes_roles/rol2_eda.md;
% no agrega interpretación nueva. Cifras: outputs/tables/eda_0*.csv.
\begin{figure}[H]
  \centering
  \figura{figuras/figura1.pdf}{\linewidth}
  \caption{(a) Distribución de la LGD por tramo (densidad). (b) LGD frente al monto original
  (escala logarítmica) y (c) frente al LTV, en entrenamiento, con tendencia LOWESS. (d) LGD de
  los créditos con y sin DTI por tramo ($n$ entre paréntesis). Las líneas punteadas marcan
  LGD $= 0$ y LGD $= 1$.}
\end{figure}

La LGD tiene cola hacia la derecha (media 0,541; mediana 0,522): el 1,9\,\% es negativa y el
8,3\,\%, mayor que 1. La validación es más dispersa que el entrenamiento (desviación estándar
0,40 frente a 0,30) y tiene más LGD sobre 1 (14,9\,\% frente a 6,6\,\%), por lo que se espera
un error mayor fuera de muestra (Figura~1a). Las decisiones se tomaron solo con entrenamiento.
La LGD baja con el monto original, sobre todo en los montos bajos, que además son los más
variables; se usa log(monto), porque su relación con la LGD es la más cercana a una recta
(Figura~1b; $r = -0{,}36$ frente a $-0{,}31$ en escala original) y es más interpretable que
Box-Cox ($\lambda = 0{,}35$). La relación con el LTV no es monótona (Figura~1c), así que ninguna
transformación la endereza, y el resto de los predictores queda en escala original. El LTV está
fuertemente correlacionado con el CLTV ($r = 0{,}85$) y con el seguro hipotecario ($r = 0{,}77$).
El seguro vale 0 en dos tercios de los créditos, por lo que entra como una indicadora de tener
seguro más el porcentaje de cobertura. Entre los factores, la LGD mediana es mayor en créditos
de inversión, refinanciamientos con retiro de efectivo y viviendas de 2 a 4 unidades. Las 21 LGD
atípicas (0,7\,\%) son pérdidas reales y se mantienen. Los 81 créditos sin DTI tienen una LGD
similar al resto (Figura~1d; Welch y Mann-Whitney, $p \geq 0{,}75$), aunque podrían
corresponder a deudores con DTI mayor que 65\,\%.
